\section{Tangent geometry and the exact nonlinear remainder}\label{sec:tangent}
\subsection{Rank-one tangent of the unchanged registration}
Let $v=(1,1,1)$, $e=(1,-1,0)/\sqrt2$, $u=a-v$, and $[v]_\times q=v\times q$. Direct multiplication gives
\begin{equation}\label{eq:tangent}% RR10
\begin{gathered}
 D=K[v]_\times
 =\frac G{\sqrt7}
 \begin{pmatrix}
 \sqrt2-\sqrt3&-\sqrt2+\sqrt3&0\\
 \sqrt2+\sqrt3&-\sqrt2-\sqrt3&0
 \end{pmatrix}
 =\ell e^T,\\
 \ell=De=\frac G{\sqrt7}
 \binom{2-\sqrt6}{2+\sqrt6},\qquad
 |\ell|^2=\frac{20G^2}{7},\qquad Ke=0,\\
 \omega_N^2=\frac{\eta^2m|\ell|^2}{2\alpha}=\frac56\nu^2 .
\end{gathered}
\end{equation}
The tangent drive $D$ has rank one; the exact $K$ has rank two. In particular, the full drive remains $f=Db+K(u\times b)$.

\subsection{A weighted forced estimate}\label{sec:forced}
For $Y=(u,b,p,r_c,z)$, introduce the Hilbert norm
\begin{equation}\label{eq:auxnorm}% RR11
 \norm Y_*^2=2\alpha(|u|^2+|b|^2)
       +\norm p_2^2+\norm{r_c}_2^2+\beta|z|^2 .
\end{equation}
This is an auxiliary proof norm, not the original native store or a newly conserved nonlinear energy. The finite shift preserves the strong regularity already supplied by P1.

On the domain with periodic $H^1$ carrier coordinates and unrestricted finite coordinates, define
\begin{equation}\label{eq:operator}% RR12
 \Ac Y=\left(
 0,-\mu D^Ty,\,
 -c(r_c)_s+\eta\B Db-\kappa\B z,\,
 -cp_s,\frac{\kappa}{\beta}y
 \right).
\end{equation}
Since $D^Ty=w\times v$, the weighted tangent native term is
$-2\alpha\mu(Db)^Ty=-\eta(Db)^Ty$. The complete cancellation is
\begin{equation}\label{eq:skew}% RR13
\begin{aligned}
 \langle Y,\Ac Y\rangle_*
 ={}&-\eta(Db)^Ty+\eta y^TDb\\
 &-c\int[p(r_c)_s+r_cp_s]\dd s
 -\kappa y^Tz+\kappa z^Ty=0 .
\end{aligned}
\end{equation}
The periodic $H^1$ product rule removes the spatial integral. The native weight $2\alpha$ is essential: using $\alpha$ would leave an uncancelled tangent term.

For an already strong curve $\mathcal Z$ satisfying
$\dot{\mathcal Z}=\Ac\mathcal Z+F$, differentiate
$(\norm{\mathcal Z}_*^2+\delta^2)^{1/2}$ with $\delta>0$.
Equation~\eqref{eq:skew} and Cauchy--Schwarz bound its derivative above by $\norm F_*$. Integrating and letting $\delta\downarrow0$ gives
\begin{equation}\label{eq:forced}% RR14
 \norm{\mathcal Z(t)}_*
 \le\norm{\mathcal Z(0)}_*+\int_0^t\norm{F(s)}_*\dd s,\qquad t\ge0 .
\end{equation}
This includes zero norm without dividing by it. Existence is inherited before this calculation, not inferred from skew cancellation.

\subsection{Complete nonlinear decomposition}\label{sec:nonlinear}
Using~\eqref{eq:nativesigns} and $f=Db+K(u\times b)$ gives the exact equation
\begin{equation}\label{eq:remainder}% RR15
\begin{gathered}
 \dot Y=\Ac Y+\mathcal N(Y),\\
 \mathcal N(Y)=\left(
 \mu w\times b,-\mu w\times u,\eta\B K(u\times b),0,0
 \right).
\end{gathered}
\end{equation}
No quadratic term is discarded. Put
$X_N=\sqrt{2\alpha}\sqrt{|u|^2+|b|^2}$ and
$V_{\rm ext}^2=\norm p_2^2+\norm{r_c}_2^2+\beta|z|^2$.
Since $|w|\le k\sqrt m\,V_{\rm ext}$, the weighted native remainder is at most $LX_NV_{\rm ext}$, while the external remainder is at most
$\eta k\sqrt m\,|u||b|\le(L/2)X_N^2$, where
\begin{equation}\label{eq:nonlinearbound}% RR16
\begin{gathered}
 L=\frac{\eta k\sqrt m}{2\alpha},\\
 \norm{\mathcal N(Y)}_*^2
 \le L^2(X_N^2V_{\rm ext}^2+X_N^4/4)
 \le L^2(X_N^2+V_{\rm ext}^2)^2,\\
 \norm{\mathcal N(Y)}_*\le L\norm Y_*^2 .
\end{gathered}
\end{equation}
The blocks are orthogonal in~\eqref{eq:auxnorm}. The second squared inequality has gap
$L^2(3X_N^4/4+X_N^2V_{\rm ext}^2+V_{\rm ext}^4)\ge0$, proving the stated global bound.
